% Body fragment for the cumulative research report. Use the main TeX file.
\newpage
\heading{13C. Encoding a grid in a permanent}\label{sec:perm}
Take finite incidence labels $I$ and surjections $v:I\to V$ and $w:I\to K$.
An incidence $n$ has a table $W_n(a,b)$ over a commutative ring $R$.
One bit is shared by all incidences in each fibre of $v$ and one by all
incidences in each fibre of $w$. The polynomial or ring element to encode is
\[
  Z(W)=\sum_{x:V\to\{0,1\}}\sum_{r:K\to\{0,1\}}
           \prod_{n\in I}W_n(x_{v(n)},r_{w(n)}).
\]
Assume $2h=1$ in $R$. E.1 gives a matrix of every order $S\ge3|I|$
whose permanent equals $Z(W)$. To see what forces the shared bits, read
each matrix row as a linear form in formal column markers. Then
\[
 \operatorname{per}(M)=
 [t_1\cdots t_S]\prod_{i=1}^S\Bigl(\sum_{j=1}^SM_{ij}t_j\Bigr).
\]
The indicated coefficient uses each column exactly once. During
multiplication we can set every $t_j^2$ to zero: any monomial with a
repeated marker already misses this coefficient, whatever the later rows
contribute.

\subhead{The first two rows choose one bit per left fibre}
Choose permutations $\sigma,\rho$ on $I$ whose cycles are respectively
the fibres of $v,w$. At each $n$ introduce three columns $A_n,B_n,C_n$.
With $A=A_n$, $B=B_n$, $C=C_n$ and $A'=A_{\sigma n}$, the first two row
forms are $h(A+B+A'-C)$ and $A+B-A'+C$. In the quotient by all
column-marker squares their product is
\[
 h\bigl((A+B)^2-(A'-C)^2\bigr)
       =2h(AB+A'C)=AB+A'C .
\]
The difference of squares cancels the mixed choices; the factor $2h$
gives coefficient one to each surviving choice. This calculation also
holds when $\sigma n=n$. Set $s_n=1$ for the term $A_nB_n$ and $s_n=0$
for $A_{\sigma n}C_n$. Column $A_m$ then occurs
\[
  \mathbf{1}_{s_m=1}+\mathbf{1}_{s_{\sigma^{-1}m}=0}
\]
times. Requiring exactly one occurrence gives $s_m=s_{\sigma^{-1}m}$
for every $m$. The remaining column at $m$ is $C_m$ if $s_m=1$ and
$B_m$ if $s_m=0$. We have recovered one independent bit for each left
fibre, and the label of its unfilled column records its value.

\smallref{\paper{Proposition E.1, pp.~31--32}.
\pub{Nc/PermCoeff.lean}{95}{Coefficient identity};
\pub{Nc/PermCoeff.lean}{143}{square-zero calculation};
\pub{Nc/PermGrid.lean}{122}{matrix entries}.}

\newpage
\heading{Recovering the weight and its sign}
Write $D_m(1)=C_m$ and $D_m(0)=B_m$. The third row at $n$ is
\[
 \sum_{a=0}^1 W_n(a,1)D_n(a) +
 \sum_{a=0}^1 W_{\rho n}(a,0)D_{\rho n}(a).
\]
Choose the first sum when $b_n=1$ and the second when $b_n=0$. The row
then uses the column at destination $m=n$ or $m=\rho n$, respectively.
Exactly one column is vacant there, forcing $a=s_m$. A destination $m$
has exactly one incoming row precisely when
$\mathbf{1}_{b_m=1}+\mathbf{1}_{b_{\rho^{-1}m}=0}=1$.
Thus $b$ is constant along each right fibre. Each destination occurs
once, and $b_n=b_m$, so the product of chosen coefficients is
$\prod_m W_m(s_m,b_m)$. Summing over the two families of bits gives
$Z(W)$. A block diagonal identity matrix adjoins as many forced
unit choices as needed to reach order $S$.

\subhead{Where the hypothesis on two enters}
For a single incidence, both permutations fix the one site. In column
order $A,B,C$ the combined coefficients in the actual matrix are
\[
 \begin{pmatrix}
     2h & h & -h\\
     0  & 1 & 1\\
     0  & W(0,1)+W(0,0) & W(1,1)+W(1,0)
 \end{pmatrix}.
\]
Its permanent is $2h\sum_{a,b}W(a,b)$. For $W(0,0)=1$ and the other
entries zero, the desired value is $1$. Choosing $h=1$ gives $2$ over
the rationals. In characteristic two, the two rows at each site are
proportional for every $h$, since both signs coincide. The permanent
there equals the determinant, so every nonempty matrix built by this
gadget has permanent zero. This identifies why E.1 does not transfer
the permanent claim to characteristic two. M.2 supplies a different
construction for Hamiltonian cycles in that characteristic.

\subhead{Raw permutations can cancel}
Take two incidences sharing both bits and tables with rows
$(1,2),(3,4)$ and $(5,6),(7,8)$. The four assignment products sum to
$5+12+21+32=70$. Treating both sites as independent would give
$10\cdot26=260$. In the actual $6\times6$ matrix the positive terms of
the permutation expansion sum to $213$ and the negative terms to $-143$.
The four legal choices emerge after collecting terms. With three sites
the program also checks that moving the shifted weight to the wrong
table changes $2924$ to $2860$. A finite-difference check also compares
every formal weight coefficient on six specified grids over the rationals.
All 512 agree, including terms that repeat an incidence.

\smallref{\pub{Nc/PermGrid.lean}{67}{Destination test};
\pub{Nc/PermGrid.lean}{175}{weight reindexing};
\pub{Nc/PermPad.lean}{68}{padding}.
\currentreview{checks/pass13c_permanent.py}{Exact checks and tables}.}

\newpage
\heading{What the circuit transfer costs in E.2--E.3}
Now take $R=F[Y]$, a polynomial ring over a field of characteristic
different from two. For each right fibre $k$, the affine form of E.1
requires that all its weights belong to the span of $1$ and a single
coordinate $Y_{o(k)}$, or are all constants. Every sum in a weighted
matrix entry uses tables from the same fibre, because $w(\rho n)=w(n)$.
The entry is therefore a constant or $aY_i+b$. The common coordinate
matters: in the singleton matrix on the preceding page, taking
$W(0,1)=Y$ and $W(0,0)=Z$ puts $Y+Z$ into a single entry.

Writing $E_{ij}$ for these entries, the identity of formal polynomials is
$\Perm_S(X)\big|_{X_{ij}=E_{ij}}=Z(W)$. If $s$ gates compute the generic
permanent, five gates suffice to replace one affine input, namely
$a,Y_i,aY_i,b,aY_i+b$. Old additions and multiplications use the outputs
of the substituted predecessor blocks; a repeated reference uses the
same stored value. The bound is $5s$ gates for $Z$. A lower bound $L$ on
the grid circuit consequently implies $s\ge L/5$.
The matrix order $3|I|$ and the circuit budget $5s$ measure different
objects. The projected polynomial has degree at most $|I|$ since each
assignment contributes a product of $|I|$ affine weights. A repeated
coordinate can occur with higher powers.

\subhead{From literals to the network polynomial}
In E.2 a nonempty clause of weight $u$ contributes, at fixed bit assignment,
\[
   1+(u-1)\prod_{j\text{ in the clause}}
                    \mathbf{1}_{\text{literal }j\text{ holds}} .
\]
The additional clause bit selects one of these two terms. The factor
$u-1$ is placed on the first literal, so it is charged exactly once.
For a complete $k$-node network with $\ell$ bits on each edge, E.3
defines a clause at each node for every possible local bit pattern.
Exactly one clause at each node contributes its variable. There are
$k2^{(k-1)\ell}$ clauses, each of length $(k-1)\ell$, and the resulting
network polynomial has degree $k$. Thus $q=k(k-1)\ell\,2^{(k-1)\ell}$
incidences suffice; E.3 uses the convenient bound
$S\ge3k^3\ell\,2^{k\ell}$. At $k=2,\ell=1$, the four literals yield
$X_{1,0}X_{2,0}+X_{1,1}X_{2,1}$ by a matrix of order 12. The looser
E.3 threshold is 96 for the same example.

\subhead{Scope of this step}
E.1--E.3 give a circuit for the network from any candidate permanent
circuit within the budget. The lower bound for that network is another
theorem. Lean rebuilt 12 publisher modules for E.1, E.2 and affine substitution.
Ten declaration types and seven axiom lists were read back. The E.3 caller
above is traced directly in its source.

\smallref{\pub{Paper/Grids.lean}{59}{Affine hypothesis};
\pub{Paper/Grids.lean}{83}{clause construction};
\pub{Nc/PermValiant.lean}{70}{gate substitution};
\pub{Paper/Flat.lean}{132}{E.3 caller}.
\currentreview{passes/pass13C-permanent.md}{Derivation and reproduction instructions}.}
